\subsubsection{روایت زمانی کامل}

ترتیب اجرای پروژه بخشی از روش‌شناسی ارزیابی است، چون اعتبار هر عدد به نسخه‌ی تقسیم و به زمانی بستگی دارد که مجموعه‌ی آزمون دیده شده است. آن مرز در این پروژه یک بار و پیش از هر مقایسه‌ای مشخص شد.

کار با سه منبع تصویربرداری آغاز شد. پس از پاک‌سازی، مواردی که نشانگر فراکتالی نامعتبر داشتند کنار گذاشته شدند و جمعیت کانونی تثبیت شد. خط اندازه‌گیری عروقی پیش از قفل تقسیم بازسازی شد و پنج نشانگر نهایی با مخرج پوشش بر پایه‌ی مساحت محتوا تعریف گردیدند. سپس تقسیم گروه‌محور ساخته و با اثر انگشت قفل شد و از آن لحظه هر مدل روی همان تقسیم اجرا شد، بدون هیچ گروه مشترک میان بخش‌ها؛ شمار هر بخش در \autoref{tab:cohort} آمده است. مقایسه‌ی اصلی یک بار روی آزمون قفل‌شده انجام شد و انتخاب مدل، آستانه و پیکربندی تنها با اعتبارسنجی صورت گرفت.

آزمایش‌های بعدی به ترکیب یادگیرنده، انتقال به منبع دیده‌نشده، ناوردایی دامنه و تکرار پرسش عروقی با یک ستون فقرات تصویری متفاوت پرداختند. در فولدهای منبع‌کنارگذاشته، منبع هدف نه در برازش، نه در پیش‌پردازش و نه در انتخاب مدل وارد نشد؛ در آزمایش ناوردایی دامنه هم قابلیت پیش‌بینی هویت دامنه از بازنمایی اندازه‌گیری شد. تحلیل‌های بعدی همه بازتحلیل اکتشافی ثانویه روی همان تقسیم قفل‌شده‌اند و مجموعه‌ی نگه‌داشته‌شده‌ی تازه‌ای برای تأیید نمی‌سازند.
